\section{Introduction}\label{sec:introduction}

We prove a uniform bound on the number of excess electrons that fixed
Coulomb nuclei can strictly bind. If there are $M$ nuclei of total charge
$Z$, that number is at most $CM$, with a constant independent of their
charges and positions. For a neutral atom with integer nuclear charge, we
also prove that the cost of removing one electron and the radius containing
all but one half of an electron each stay between positive universal
constants. These conclusions address the number of bound electrons, the
first ionization energy, and the location of the outer electronic mass.
Their proofs share the local screening estimates developed in this paper.

\subsection{The model and the results}

Fix $M\ge1$, distinct positions $R_1,\ldots,R_M\in\R^3$, and real
charges $z_1,\ldots,z_M\ge1$. Put
\[
 Z=\sum_{j=1}^M z_j,\qquad \nu=\sum_{j=1}^M z_j\delta_{R_j},
 \qquad K(x)=|x|^{-1},\qquad V=K*\nu.
\]
For $n\ge1$ let $\mathcal H_n=\bigwedge^nL^2(\R^3;\C^2)$, so that
there are two electron spin states and antisymmetry exchanges spatial and
spin variables together. The Hamiltonian is the self-adjoint operator
associated with the closed, bounded-below Coulomb form on
$\mathcal Q_n=\mathcal H_n\cap H^1((\R^3)^n;(\C^2)^{\otimes n})$:
\begin{equation}\label{eq:hamiltonian}
 H_n=\sum_{i=1}^n\left(-\frac12\Delta_i-V(x_i)\right)
       +\sum_{1\le i<j\le n}|x_i-x_j|^{-1}.
\end{equation}
The nuclei are fixed; their mutual repulsion is a scalar independent of
$n$ and is omitted. Write $E_n=\inf\sigma(H_n)$ and $E_0=0$.
We say that sector $n$ \emph{binds strictly} if $E_n<E_{n-1}$.

\begin{theorem}[Uniform excess charge]\label{thm:main}
There is a universal finite constant $C$ such that every nuclear system
specified above and every integer $n\ge1$ satisfy
\[
 E_n<E_{n-1}\quad\Longrightarrow\quad n\le Z+CM.
\]
Moreover, $E_n=E_{n-1}$ whenever $n>Z+CM$. For one nucleus this gives
$n\le Z+C$ under strict binding, for every real $Z\ge1$.
The same constant is independent of all charges and nuclear positions.
\end{theorem}

The positions may be arbitrarily close or far apart. The theorem resolves
positively the excess-charge form of the ionization conjecture in this
strict-binding formulation \citep{Simon2000,Lewin2025OpenProblems}.
Equality of adjacent energies does not by
itself exclude a ground state at threshold; such threshold states are not
needed here.

For the radius conclusion specialize to one nucleus at the origin with
integer charge $Z\ge1$. Write $H_{n,Z}$ and $E(n,Z)$ when the nuclear
charge needs to be displayed. If $\Psi_Z$ is a normalized ground state of
$H_{Z,Z}$, its spin-summed density is
\[
 \rho_{\Psi_Z}(x)=Z\sum_{\sigma_1,\ldots,\sigma_Z}
 \int |\Psi_Z(x,\sigma_1;x_2,\sigma_2;\ldots;x_Z,\sigma_Z)|^2
                         \,dx_2\cdots dx_Z,
\]
where the integral is absent for $Z=1$. Its mass is $Z$. Define
\begin{equation}\label{eq:radius}
 R(\Psi_Z)=\inf\left\{r\ge0:
       \int_{|x|>r}\rho_{\Psi_Z}(x)\,dx\le\frac12\right\}.
\end{equation}

\begin{theorem}[Uniform outer half-electron radius]\label{thm:radius}
There are universal constants $0<c\le C<\infty$ such that, for every
integer $Z\ge1$ and every normalized ground state of $H_{Z,Z}$,
\[
 c\le R(\Psi_Z)\le C.
\]
\end{theorem}

The existence of these ground states is included in
Lemma~\ref{lem:atomic-binding}. No uniqueness or spherical symmetry is
assumed. The exterior mass in \eqref{eq:radius} stays fixed as $Z$ grows;
it describes outer electrons rather than a fixed fraction of the atom.

\begin{corollary}[Two-sided first-ionization bound]\label{cor:first-ionization}
For the full nonrelativistic Coulomb atom specified above, with integer
nuclear charge $Z\ge1$, one-electron kinetic term $-\Delta/2$, and two
electron spin states, write $E_Z(n)=E(n,Z)$. There are universal constants
$0<c\le C<\infty$ such that the first ionization energy of the neutral atom
satisfies
\[
 c\le I_1(Z):=E_Z(Z-1)-E_Z(Z)\le C
 \qquad\text{for every integer }Z\ge1.
\]
\end{corollary}

\subsection{History and the scales of the problem}

Simon's atomic formulation asks whether $N_0(Z)-Z$ stays bounded, where
$N_0(Z)$ is the first sector at which the energy equals every later sector
energy \citep[Problem~9]{Simon2000}. Since the sector energies are
nonincreasing and ultimately constant, this is the largest strictly
binding sector used here. The corresponding molecular conjecture has
the form $N_c\le Z+CM$ \citep[Section~2.1.2]{Lewin2025OpenProblems}.

The early charge bounds already distinguished atoms from molecules.
Lieb proved bounds linear in the total nuclear charge for fixed atoms and
molecules, without requiring Fermi statistics \citep{Lieb1984}.
For fermionic atoms, Lieb, Sigal, Simon and Thirring established asymptotic
neutrality: the maximal strictly bound particle number $N_c(Z)$ satisfies
$N_c(Z)/Z\to1$ as $Z\to\infty$ \citep{LSST1984,LSST1988}.
Fefferman and Seco, and Seco, Sigal and Solovej, subsequently obtained
quantitative sublinear bounds \citep{FeffermanSeco1990,SecoSigalSolovej1990}.
Ivrii obtained further semiclassical charge bounds for atoms and heavy
molecules under the hypothesis that each nuclear charge is comparable
to the electron number, with improved exponents for atoms and for
suitably separated nuclei \citep[Theorem~5.2.1]{Ivrii2012}.
For explicit finite-charge estimates, Nam proved
$N_c(Z)<1.22Z+3Z^{1/3}$ \citep{Nam2012}.
More recently, Hundertmark, Pattakos and Schulz obtained
$N_c(Z)<1.1185Z+4Z^{1/3}$ for $Z\ge4$ \citep{HPS2025}.
These estimates supply useful quantitative control while leaving the
uniform excess-charge assertion distinct.

For the full many-electron Schr\"odinger atom, Seco, Sigal and Solovej
also obtained, for large $Z$, an upper bound for the neutral first ionization energy
and a lower bound for an exterior-mass radius, both given by powers of
the nuclear charge \citep[Theorems~1 and~3]{SecoSigalSolovej1990}.
Their radius leaves one expected electron outside, whereas
\eqref{eq:radius} leaves one half. Their energy upper bound grows with
$Z$, and their radius lower bound decays with $Z$.
The upper bound in Corollary~\ref{cor:first-ionization} establishes the
usual boundedness assertion for $I_1(Z)$
\citep[Section~2]{Solovej2016}; the uniform positive lower bound is an
additional conclusion.

The statistical theories introduced by Thomas and Fermi
\citep{Thomas1927,Fermi1927} predict a particularly rigid exterior
screening profile. Lieb and Simon established the variational theory and
its leading large-charge energy and density limits
\citep{LiebSimon1973,LiebSimon1977}. Uniform ionization bounds were
proved in Thomas--Fermi--von Weizs\"acker theory by Benguria and Lieb
\citep{BenguriaLieb1985}, and in reduced and unrestricted Hartree--Fock
theory by Solovej \citep{Solovej1991,Solovej2003}. Solovej's unrestricted
Hartree--Fock analysis compares screened potentials through successive
spatial scales and also establishes uniform bounds and asymptotics for
outer-electron radii \citep[Definition~1.2, Theorem~1.5, and Sections~10--13]{Solovej2003}.
Those conclusions concern their stated variational models; the states in
\eqref{eq:hamiltonian} may have arbitrary many-body correlations.

The bulk density limit lives on the length
scale $Z^{-1/3}$, whereas the normalized mass outside the radius in
\eqref{eq:radius} is $1/(2Z)$. Its vanishing fraction makes bulk convergence
insufficient to prove Theorem~\ref{thm:radius}. The fixed exterior-mass
viewpoint is also central to the generalized ionization conjecture
\citep[Section~3]{Solovej2016}.

Companion papers use the estimates developed here, with additional arguments,
to obtain Thomas--Fermi asymptotics for full two-spin Coulomb atoms.
The energy result sends the number of removed electrons to infinity while
it remains negligible relative to the nuclear charge
\citep[Theorem~1.1]{OpenAI2026GeneralizedEnergy}; the neutral-radius result
first takes upper and lower large-charge limits at fixed expected exterior
electron mass, and then lets that mass tend to infinity
\citep[Theorem~1.1]{OpenAI2026GeneralizedRadius}.

The analytic ingredients have separate roles. The Lieb--Thirring kinetic
inequality provides a positive coarse density bound, for which we use
Rumin's spectral-splitting argument
\citep{LiebThirring1975,Rumin2011}. The exact Thomas--Fermi coefficient
comes instead from coherent-state filling and Neumann eigenvalue counting,
with explicit remainders. These constructions have classical antecedents
in the semiclassical analysis of Lieb and Simon and in the coherent-state
comparison of Solovej
\citep[Theorems~III.10--III.13]{LiebSimon1977}
\citep[Lemma~8.2]{Solovej2003}. Spatial square partitions are the usual
IMS localization; sector-valued localization for correlated states is
systematically developed by Lewin \citep[Proposition~3.1 and Example~3.3]{Lewin2011Geometric}.
Our conditional core construction retains the removed positions and spins,
so that the remaining correlated state can be compared with local trial
electrons through its own screened field.

\subsection{Obstacles and the common argument}

Fixing the total electron number obstructs a local density comparison:
a ball may contain too many or too few electrons, and its conditional core
need not minimize any fixed local-particle-number problem. We therefore
first work relative to the minimum $E=\inf_{n\ge0}E_n$ over all sectors.
A coarse binding bound guarantees that it is attained at a largest
strictly bound sector. A normalized state with form energy at most
$E+\delta$ will be called a $\delta$-state. Keeping $\delta$ visible makes
local estimates available again after an event is imposed, after particles
are deleted, and after a shell weight changes the probability law.

Three steps are common to the molecular and atomic conclusions. First, a
spatial cut records every removed electron and leaves an antisymmetric
conditional core. In the hole, arbitrary expected trial mass is allowed
because each actual trial has an integer sector with energy at least $E$.
The coherent upper and Neumann lower comparisons have the same kinetic
coefficient. Their comparison produces a local Thomas--Fermi minimizer
$\rho$ and its screened field $W$, related by
\[
 \rho=k(W_+)^{3/2},\qquad k=\frac{2^{3/2}}{3\pi^2}.
\]
It also controls the Coulomb discrepancy from the actual fine density
and the energy of that density in the negative part of $W$. A bootstrap
for the positive field closes a bound on the second moment of local
electron counts.

Next, smoothing and noisy observations produce a conditional density
$\mu$ and a screened field. In the atomic case the field is
$H=V-K*\mu$. On nucleus-free patches we transfer the local Thomas--Fermi
law to these quantities, as one-sided comparisons between $\mu$ and
$k(H_+)^{3/2}$ at positive, bounded rescaled heights. The molecular
comparison retains the nearby nuclear potential explicitly: after
splitting off its singular part, we compare the continuous conditional
offset through an averaged Thomas--Fermi response.

To prove these comparisons, observe all positions with independent errors
of width $\ell$ and forget the particle labels. An observation event of
probability $p$ has a symmetric square-root likelihood multiplier.
Applied to a sector eigenstate, its energy cost is bounded by
$C\ell^{-2}\log^5(e/p)$, without a particle-number factor. A hypothetical
failure can therefore be tested by the local quantum estimates. The
original observation posterior and the fresh patch posterior are
different: two separate conditional-expectation identities compare their
fields after averaging in the event law. The negative-field penalty
controls the rare configurations needed to pass a lower field estimate
to expectation.

Finally, a positive comparison field is continued outward. A local
maximum chooses between a shifted nonlinear subsolution and a shifted
observed field. The inverse comparisons let us replace the model reaction
by the actual conditional density on one more layer. Conditional averaging
then forgets one observation array, with convexity preserving the nonlinear
inequality. The errors are summed through their expected integrals. At the
terminal scale, comparison with a compact source potential and its spherical
mean turns the positive field into a lower bound on the nuclear charge
minus the enclosed electron mass, up to the accumulated error.

The geometry determines how this mechanism closes. For a molecule, the
local length measures nearby nuclear charge. Empty nuclear annuli must be
counted without a loss for their number of scales. A decomposition by
nested nuclear representatives gives a total cost proportional to $M$;
it bounds the electrons outside the terminal region and closes the
excess-charge estimate.

For an atom, local balls away from the nucleus use their distance from it;
this distinct scale reaches inside the bulk length. The continuation
produces an interior electron deficit for sector ground states with
$n\le3Z$ and bounded energy offset from $E$. We first use it at offset zero. Deleting exterior particles then
proves $E_{Z-1}\le E+C$, so both neutral and singly ionized ground states
become eligible for the fixed-offset estimates. The neutral deficit gives
the lower radius bound. A spatially averaged insertion into the singly
ionized state gives a uniform positive gap $E_{Z-1}-E_Z\ge c$. Together
with the offset bound this proves Corollary~\ref{cor:first-ionization},
after adjusting the constants at finitely many smaller charges. The gap
also yields the expected tail bound $C/R$ through a shell-weighted change
of law. Here the shell weight is measurable in the revealed data, which
is why this change of law preserves the surviving conditional field.

\subsection{Organization and normalization}

Section~\ref{sec:foundations} provides the sector facts, and
Sections~\ref{sec:localization}--\ref{sec:observations} prove the shared
localization, screening, and conditioning estimates. The two local
geometries are specified in Section~\ref{sec:screening}. From these shared
tools, Sections~\ref{m:sec:void}--\ref{m:sec:outward} prove the molecular
excess-charge bound. A reader interested in the neutral atomic results
may instead proceed directly to Sections~\ref{a:sec:conditioning}--\ref{a:sec:tail};
these proofs do not use the molecular excess-charge theorem. The atomic
comparison and intermediate subsolutions come first, followed by deletion,
insertion, and the tail argument in the order described above.
Section~\ref{sec:interfaces} collects the precise inputs for the companion
asymptotic problems, including the classical Thomas--Fermi comparisons
proved in Section~\ref{m:sec:barriers}.

All density integrals include spin summation. For a real number or
function $v$, write $v_+=\max(v,0)$ and $v_-=\max(-v,0)$. The constants
$C,c>0$ may change from line to line and are universal unless a dependence
is stated. With the kinetic normalization in \eqref{eq:hamiltonian},
\begin{equation}\label{eq:tf-constants}
 c_{\rm TF}=\frac3{10}(3\pi^2)^{2/3},\qquad
 k=(5c_{\rm TF}/3)^{-3/2}=\frac{2^{3/2}}{3\pi^2},\qquad c_F=4\pi k.
\end{equation}
The coarse kinetic lower bound does not assert that its unspecified
constant equals $c_{\rm TF}$.
